Phase~5 runs a TTM-inspired trace-distance proxy: at each step it accumulates the trace
distance between the actual trajectory and the one-step Markovian channel prediction. It
is a bounded backflow diagnostic, not a transfer-tensor reconstruction. The naming is
deliberately comparative: the full Transfer Tensor Method reconstructs non-Markovian
memory through transfer tensors inferred from process data~\cite{CerrilloCao2014},
whereas QSOT2 uses only a trace-distance proxy in that spirit.

Two trajectories are kept distinct:

\begin{itemize}[leftmargin=1.5em]
\item \emph{Synthetic self-test.} An artificial coherence backflow is injected to confirm
  the detector registers $\mathrm{nm} > 0$ (reference value
  $\mathrm{nm} \approx 1.4\times10^{-4}$ at $\alpha=0.1$). A pass validates the detector
  only.
\item \emph{Model-generated trajectory.} Produced by applying the configured channel
  consistently with no injected revival, this yields $\mathrm{nm} = 0$: the
  phenomenological map is Markovian by construction.
\end{itemize}

The detection threshold and the Ricci norm at which it was evaluated are surfaced with
the result, so the (here non-engaged) curvature scaling is explicit rather than implied.
A zero model-trajectory result is a property of the implemented map, not a deep physical
discovery.
